\section{سؤالی با چند بخش \points{۱۲}}

هر \lr{\texttt{section}} یک سؤال است: یک نوار ضخیم، بعد شماره‌ی سؤال در
حاشیه با اندازه‌ی درشت، عنوان کنارش و بارم در کادری تیره در انتهای سطر. متن
سؤال به‌شکل معمولی از اینجا ادامه پیدا می‌کند.

\marginlabel{یادداشت‌های کناری اینجا می‌نشینند: راهنمایی، تقسیم بارم، یا
ارجاع به اسلایدی که سؤال از آن آمده است.}

\subsection{بخش نخست}

زیربخش‌ها همان قالب را یک اندازه کوچک‌تر می‌گیرند: شماره در حاشیه، یک خط نازک
زیر عنوان، بدون نوار ضخیم.

\subsection{بخش دوم}

فهرست‌های شماره‌دار، شماره‌ی خانگی قالب را می‌گیرند که از دل مربعی تیره
بیرون آمده است:

\begin{enumerate}
    \item ناوردا را بنویسید.
    \item نشان دهید در گام نخست برقرار است.
        \begin{enumerate}
            \item حالت پایه.
            \item گام استقرا.
        \end{enumerate}
    \item نتیجه بگیرید.
\end{enumerate}

شماره‌هایی مثل \hwnum{4.1} در متن راست‌به‌چپ وارونه می‌شوند؛ آن‌ها را با
دستور \lr{\texttt{hwnum}} بنویسید.

\section{سؤالی بدون زیربخش \points{۸}}

بعضی سؤال‌ها یک پاراگراف‌اند و یک خواسته. آن‌ها به ساختار نیازی ندارند؛ همان
عنوان و صورت سؤال کافی است.

\begin{asidebox}{راهنمایی}
    از تعریف شروع کنید، نه از نتیجه‌ای که می‌خواهید به آن برسید.
\end{asidebox}
